\chapter{Background and Related Work}\label{chap:background}

%TL;DR: The chapter surveys the concepts needed before the system chapters.
Chapter~\ref{chap:intro} provided motivation for scene-aware dialogue for the Pepper robot. This chapter makes that motivation technically precise. It does not yet describe the implemented architecture in detail.\footnote{This is described in detail in the following Chapters~\ref{chap:system_architecture_rewrite}, \ref{chap:server_side_rewrite}, and~\ref{chap:robot_side_rewrite}} Instead, it formally introduces the terms, tasks, and methods as well as the related work needed to understand the following implementation chapters.

%TL;DR: The chapter follows the conceptual path from HRI to perception, memory, and dialogue.
In Section~\ref{sec:hri-social-robots} the thesis project is redefined in the context of HRI and social robotics. Section~\ref{sec:problem-setting-scene-aware-dialogue} hints at the implementation from more formal point of view preparing grounds for further explanation of each component. The rest of sections follows the order of input processing.
%Since Pepper is not just a camera or a chatbot, but rather an ebomdied interaction partner, we first discuss human-robot interaction (HRI) and social robots. MIGHTBEDROPPED: The chapter then formalizes the problem of scene-aware dialogue, identifying the key system requirements that motivate the components described in the following sections.

It starts with object detection, followed by tracking, and re-\-iden\-ti\-fi\-ca\-tion, which provide object-level observations and persistent identities over time. It then discusses grounding and scene graph generation, where detections become structured relational representations. The final technical sections introduce language models, vision-language models, retrieval and caching methods, and the Pepper platform.

\section{Human-Robot Interaction and Social Robots}\label{sec:hri-social-robots}

%TL;DR: HRI is introduced as the field that makes interaction quality part of the technical problem.
Human-robot interaction (HRI) studies how humans and robots interact, how these interactions can be designed, and how they should be evaluated. Goodrich and Schultz describe HRI as an interdisciplinary field in which robotics, human factors, psychology, design, linguistics, and computer science meet~\cite[203--216]{goodrich_human-robot_2007}.

%TL;DR: Goodrich and Schultz give attributes for describing HRI systems.
From a design perspective, Goodrich and Schultz characterize an HRI problem using five attributes: the level and behavior of autonomy, the nature of information exchange, the structure of the human-robot team, adaptation and learning of both the human and the robot, and the shape of the task~\cite[216--217]{goodrich_human-robot_2007}.

%TL;DR: The five attributes define autonomy, communication, roles, adaptation, and task shape.
The level of autonomy determines how independently the robot acts and how control is shared between the human and the robot. Information exchange describes properties of communication such as modality, timing, and clarity. Team structure captures how the human-robot system is organized, including the number of participants and their roles (e.g., one human operating one robot or one human supervising multiple robots). Adaptation and learning describe how the human and/or the robot adjust their behavior over time based on interaction. The shape of the task characterizes the structure of the problem being addressed, including the nature of the goals, the degree of openness, and the constraints under which the interaction takes place~\cite[216--232]{goodrich_human-robot_2007}.

%TL;DR: Sheridan complements design attributes with application domains.
While these attributes describe how HRI problems can be characterized, they do not distinguish where such systems are deployed. To extend this perspective, it is useful to consider major application areas of HRI.

%TL;DR: Sheridan's taxonomy distinguishes social interaction from other HRI areas.
Sheridan distinguishes four broad applications of HRI: supervisory control, teleoperation, automated vehicles, and social interaction~\cite[525--526]{Sheridan2016HRI}. Supervisory control involves high-level human oversight of partly autonomous systems. Teleoperation denotes direct remote control of a robot, typically in remote or hazardous environments. Automated vehicles operate largely autonomously, with humans acting as passengers, supervisors, or fallback operators. Social interaction focuses on robots whose primary function depends on direct communication and interaction with humans.

%TL;DR: Social robots arise from the social interaction domain.
The last category, that of social interaction, gives rise to the notion of social robots. The term is not used uniformly in the literature. Henschel, Laban, and Cross note that there is no universally accepted definition, but that common characteristics include embodiment, some degree of autonomy, and interactional behavior interpreted by humans as social~\cite[11--12]{henschel_what_2021}. %In this thesis, a social robot is understood as a physically embodied system that communicates or cooperates with humans in a socially interpretable manner.

%TL;DR: User studies identify traits that make robots appear social.
User studies further elaborate what makes a robot appear social in practice. De Graaf et al. identify eight characteristics influencing perceived socialness: the capability for two-way interaction, the ability to display thoughts and feelings, social awareness of the environment, provision of social support, autonomy, and, less prominently, cosiness, similarity to the user, and mutual respect~\cite{de_graaf_what_2015}. The latter three were mentioned less frequently, suggesting lower importance compared to interaction, awareness, and responsiveness.

%TL;DR: Social robots do not need to be humanoid.
Importantly, social robots are not required to resemble humans in form. Sociality is not tied to anthropomorphic appearance, but rather to how interaction is perceived by humans, which can be documented well by the use of pet-like robots~\cite[10]{henschel_what_2021}. 

%TL;DR: Social robots remain constrained by their technical nature.
At the same time, these systems remain machines, and their social qualities are constrained by their underlying technical capabilities. This tension between perceived social behavior and actual system competence is central to the design of socially interactive robots~\cite{cross_mind_2021}.


%TL;DR: Pepper is positioned as a social interaction platform.
%The present work belongs primarily to the domain of social interaction. Pepper is not used here as a high-payload manipulator, a remotely controlled machine, or an autonomous vehicle, but as a platform for proximate interaction with a human user. The central question is therefore not only whether the robot can detect objects, but whether these detections can support understandable and context-sensitive dialogue.

%TL;DR: Social embodiment raises expectations.
SOMEBRIDGINSENTENCE
%Robot embodiment changes user experience and expectations. A user may reasonably expect that a robot sharing the environment can refer to visible objects, distinguish people, and remember previously mentioned items. Henschel et al. discuss studies of Pepper users in which reciprocal conversation and the limitations of single-turn interaction become salient expectations~\cite[12]{henschel_what_2021}. Similar effects appear in work on LLM-powered HRI, where fluent language leads users to expect broader situational competence than the system actually possesses~\cite{kim_understanding_2024,atuhurra_leveraging_2024}.

Robot embodiment changes user experience and expectations. On the one hand, physical presence provides a natural form of grounding, enabling interaction to be anchored in shared space. On the other hand, embodiment raises expectations about the robot's capabilities. A user may reasonably assume that a robot sharing the environment can support richer, more reciprocal interaction.

Henschel et al. report that users interacting with the humanoid robot Pepper expected it to engage in reciprocal conversation, and were disappointed when the interaction was limited to single-turn exchanges~\cite[12]{henschel_what_2021}. This illustrates how embodiment can amplify expectations beyond what the system can deliver. Similar effects appear in work on LLM-powered HRI, where fluent language leads users to expect additional interaction capabilities that are not reliably supported by the system~\cite{kim_understanding_2024}.


Moreover, Gava et al.\ found that, despite using the same conversational framework, participants rated an embodied robot more positively than an app. They compared the two setups on perceived reliability, awareness of system behavior, speech quality, and ease of correcting mistakes~\cite{gava2024physical}. This suggests that physical embodiment is an important part of social robotics and it may lead users to judge conversational failures more favorably. Although the result should be interpreted cautiously because the study was preliminary and involved a rather small sample.

%TL;DR: Related systems motivate grounding.
%Several recent systems connect social robots with LLMs, VLMs, or external computation. Some focus on human support in teaching and education~\cite{latif_physicsassistant_2024,sievers_humanoid_2025}. Others focus on assistance, libraries, industry, or general conversational flexibility~\cite{trinquet_future_2025,brad_retrieval-augmented_2025,becerra_improving_2024}. These systems address a common challenge: while modern language models enable flexible interaction, the robot still requires grounding in physical context.

%TL;DR: The system is characterized using HRI attributes.
Finally, using the presented attributes of Goodrich and Schultz, this thesis HRI setting can be described more rigorously. Pepper has limited, interaction-triggered autonomy, exchanges information through speech, vision, tablet output, and operates in a proximate one-human-one-robot team structure. Adaptation is realized as interaction-time updates to memory and dialogue state rather than online model learning. The task is open-ended but bounded, consisting of situated dialogue about the observed environment. The system belongs to the social interaction domain of HRI.

This specification makes the HRI problem addressed in this thesis explicit; Enabling a physically embodied robot to maintain consistent, context-aware dialogue grounded in its current and accumulated perception under the constraints of limited autonomy and interaction-time state updates. We believe this framing clarifies the expectations placed on the system and provides a concrete basis for the design decisions presented in the following chapters.

%TL;DR: HRI motivates grounding across perception, memory, and dialogue.
% We can therefore summarize the HRI background as follows. If Pepper were only a camera, object detection would be sufficient. If it were only a chatbot, language modeling would be sufficient. Because it is an embodied social robot, the system must preserve a connection between what the robot says, what it sees, what it has seen before, and what the user can plausibly expect it to know.

\section{Problem Setting: Scene-Aware Dialogue}\label{sec:problem-setting-scene-aware-dialogue}
MIGHT BE OFFLOADED TO CHAPTER \ref{chap:system_architecture_rewrite} WHERE ARCHITECTURE OVERVIEW IS PROBABLY MORE FIT

%TL;DR: Scene-aware dialogue is defined as a system-level task.
Before surveying individual methods, we need to state what kind of problem they support. In this thesis, scene-aware dialogue is not simply dialogue accompanied by an image. It is dialogue in which the robot maintains a representation of its perceived environment and uses that representation when answering the user. The problem therefore combines visual perception, identity persistence, relation extraction, memory update, context selection, and language modeling.

Formally speaking, the system operates over a sequence of observations \(O_1, O_2, \ldots, O_T\), where each observation \(O_t\) consists of an image together with robot metadata (e.g., camera pose, head orientation, or other perception signals coming from NAOqi API).

From each observation, the system extracts scene information and updates an internal memory. The memory at time \(t\) can be written as
\[
M_t = (\mathcal{E}_t, \mathcal{R}_t, \mathcal{H}_t),
\]
where \(\mathcal{E}_t\) is a set of persistent entities (objects), \(\mathcal{R}_t\) represents relations between them (e.g., a scene graph), and \(\mathcal{H}_t\) stores relevant history from past observations or dialogue.

The memory is updated incrementally:
\[
M_t = U(M_{t-1}, O_t),
\]
where the update function \(U\) includes object detection, feature extraction, identity association (re-identification), relation extraction, and merging into the existing memory.

This makes explicit that the system does not answer only from the current observation, but from an accumulated representation of the scene.

On dialogue level, at turn \(k\), the user provides a query \(q_k\). The system selects relevant context \(C_k\) from memory and generates a response using a language model $p_\theta$:
\[
y_k \sim p_\theta(y \mid q_k, C_k).
\]

The main design question is how to construct \(C_k\). Broad questions (e.g., \enquote{what do you see?}) require a summary of the scene, while targeted questions (e.g., \enquote{what is next to the laptop?}) require retrieving a specific subset of entities and relations.

In summary, the system consists of the following components: processing observations (image + metadata), object detection and feature extraction, identity matching across observations, relation extraction, memory update, context selection, and response generation. The following sections introduce these components in detail.

% Offloaded to Chapter~\ref{chap:system_architecture_rewrite}: the full architecture-level formalization of memory updates, context construction, and dialogue turn flow.

\section{Object Detection, Tracking, and Re-Identification}\label{sec:object-detection-tracking-reid}

%TL;DR: Object detection provides localized object hypotheses.
The first technical requirement for scene awareness is object detection. Object detection is the task of identifying object instances in an image and localizing them, most commonly by bounding boxes. Formally, a detector maps an image \(I\) to a finite set
\[
\mathcal{D}(I)=\{(b_i,\ell_i,s_i)\}_{i=1}^{N},
\]
where \(b_i\in\mathbb{R}^4\) is a bounding box, \(\ell_i\) is a class label, and \(s_i\in[0,1]\) is a confidence score~\cite[Section 50.4.1]{foundationsCVbook}. 

A traditional object-detection pipeline can be described in three stages. First, the system proposes candidate bounding boxes. Second, each candidate region is classified, often with an additional bounding-box refinement step. Third, the resulting boxes are post-processed to remove low-confidence and redundant detections. The standard overlap measure for boxes is Intersection over Union (IoU), a ratio of the area of the intersection of two candidate bounding boxes and their union area. A more complex redundancy-removal method, non-maximum suppression (NMS) uses this overlap as well as the confidence of individual detections to keep the highest-confidence box while suppressing nearby boxes that likely refer to the same object.~\cite[Section 50.4.1.2]{foundationsCVbook}. 


Modern object detection builds on the development of convolutional neural networks (CNNs) which were the main architecture for neural networks with image data. CNNs exploit the spatial structure of data by applying learned local filters and reusing them across different image positions. This makes them well suited to visual data, where local patterns such as edges or textures appear at different locations. First CNN layers tend to represent simple local structures, while deeper layers compose these structures into more abstract features which in turn are useful for the downstream task of, say, recognition and detection~\cite[439]{lecun_deep_2015}. 

Moreover, CNNs bear an immense importance in neural network history, as Goodfellow et al.\ notes, they were among the first deep models to succeed commercially, including early systems for check reading and handwriting recognition, making the way for other neural architectures too~\cite[365]{Goodfellow-et-al-2016}. 


Early deep-learning detectors used CNNs inside the traditional proposal-based pipeline. In R-CNN, the candidate regions were not generated by the CNN itself. They were produced by selective search, an external class-agnostic proposal method that groups pixels into small segments and then hierarchically merges similar neighboring regions using low-level cues such as color, texture, size, and spatial compatibility~\cite{uijlings_selective_2013}. The resulting boxes are only hypotheses about where objects may be; they do not yet determine the object class. 

R-CNN then crops each proposed region from the image, warps it to a fixed size, applies a CNN to extract features, and classifies these features with a separate classifier~\cite{girshick_rich_2014}. This improved detection accuracy, but it was computationally expensive because the CNN had to be evaluated separately for many proposed regions. Fast R-CNN improved this by computing convolutional features once for the whole image and then pooling features for each proposed region from the shared feature map~\cite{girshick_fast_2015}. 

Faster R-CNN removed the external selective-search step by introducing a Region Proposal Network (RPN), a small fully convolutional network that slides over the shared feature map and predicts objectness scores and box offsets for predefined anchor boxes~\cite{ren_faster_2016}. These methods show the transition from hand-designed proposal pipelines toward learned detectors, but they still preserve a two-stage structure: propose regions, then classify and refine them.

Different, more up-to-date approach is taken by YOLO, which is short for \emph{You Only Look Once}. Instead of going the original path of generating many candidate regions and classifying them separately, YOLO processes the whole image in a single forward pass through one neural network~\cite{redmon_you_2016}. This is the reason behind the \enquote{one look}; The network sees the full image once, rather than repeatedly evaluating many cropped windows or region proposals as was the norm before. 

In the original YOLO formulation, the image is first divided into a spatial grid. The convolutional part of the network extracts image features, and each grid cell predicts bounding boxes, confidence scores, and class probabilities for objects whose centers fall inside that cell. In this way, localization and classification are predicted jointly, in one look; The model directly predicts object boxes and labels from the global image representation, without first generating a separate list of candidate regions. This is why YOLO can be computationally more efficient than scanning-window or proposal-based methods~\cite[Sec.~50.4.1.4]{foundationsCVbook}.

Modern YOLO models build on this efficiency, but they also differ from the original architecture while preserving the conceptual structure. The \enquote{backbone} extracts visual features from the image. The \enquote{neck} combines features across multiple scales, which helps detect both small and large objects. The detection \enquote{head} predicts bounding-box coordinates, objectness/confidence values, and class scores. The predictions are then filtered by confidence and typically post-processed with NMS. YOLO models by Ultralytics, a company specializing in computer vision, are a prominent and widely used modern implementation family, they are released in different sizes to trade accuracy against latency and compute cost. Their uses are well documented, and range from robotics to agriculture, which clearly states their practicality~\cite[1, 4, 10--11]{sapkota_ultralytics_2026}.


Next paradigm shift in object detection is transformer-based object detection. DETR, detection transformer, formulates detection as direct set prediction~\cite{carion_end--end_2020}. An image backbone first extracts visual features. A transformer encoder processes these features using self-attention, and a transformer decoder uses learned object queries to produce a fixed-size set of object predictions.\footnote{For formal definition of self-attention and a transformer see Section \ref{sec:llms-vlms}.}

This means that the model does not produce an ordered list of dense local predictions that must later be cleaned by NMS. 
Each query predicts either one object or a special no-object class. During training, predictions are matched to ground-truth objects with bipartite matching, which encourages one prediction per object and reduces the need for NMS. Instead, it predicts a fixed-size set of object hypotheses, where each hypothesis should correspond either to one object in the image or to a special \enquote{no object} class.

The difference between YOLO and DETR is therefore structural. YOLO is a dense one-stage detector: spatial locations in feature maps directly predict boxes and class scores. DETR is a set-based detector: learned queries predict a set of object hypotheses. 

% RT-DETR adapts the DETR idea for real-time detection and shows that transformer detectors can compete with YOLO-style models under practical latency constraints~\cite{zhao_detrs_2024}. CHECK: Transformers handle occlusion better: \cite{sapkota_rf-detr_2025} 

The disadvantage of the original DETR is computational cost and training difficulty. RT-DETR addresses this by keeping the end-to-end DETR formulation while making the architecture more efficient~\cite{zhao_detrs_2024}. Its hybrid encoder decouples two operations that are expensive if done naively: interaction within each feature scale and fusion across different feature scales. It also introduces uncertainty-minimal query selection, which initializes the decoder with higher-quality object queries instead of relying only on learned query slots. In addition, RT-DETR can trade accuracy for speed by reducing the number of decoder layers at inference time, without retraining.

RF-DETR, developed by Roboflow, follows the real-time DETR direction but focuses on optimizing the architecture for accuracy-latency trade-offs~\cite{robinson_rf-detr_2026}.
The main technical idea in RF-DETR is weight-sharing neural architecture search (NAS). Neural architecture search means that the architecture is treated as a search problem: rather than manually fixing all design choices, the method evaluates many possible configurations and selects those with good accuracy-latency trade-offs. In ordinary NAS, this can be very expensive because each candidate architecture may need to be trained separately. RF-DETR avoids this by training one shared model that contains many possible sub-networks~\cite{robinson_rf-detr_2026}. 

During training, RF-DETR samples different configurations from a predefined search space and updates the shared weights under those configurations. The searched choices include input image resolution, patch size, number of windowed-attention regions, number of decoder layers, and number of query tokens. In other words, the search does not simply tune scalar hyperparameters after training; it changes which parts of the model computation are active and how expensive inference will be. However, it does not train thousands of independent models. The candidate architectures share parameters during training ~\cite{robinson_rf-detr_2026}.

% RF-DETR, developed by Roboflow, follows this real-time detection-transformer direction and uses neural architecture search to find accuracy-latency trade-offs~\cite{robinson_rf-detr_2026}. WHAT IS NEURAL ARCHITECTURE SEARCH AND HOW IS IT USED!!!

YOLO, RT-DETR, and RF-DETR are normally used as closed-vocabulary detectors. Their class label \(\ell_i\) is selected from a fixed ontology learned during training or fine-tuning. COCO is the standard example of such an ontology: it contains 80 common object categories and is widely used for training and benchmarking object detectors~\cite{lin2015microsoftcococommonobjects}.

COCO-pretrained models are therefore useful general detectors, but they cannot directly predict arbitrary labels outside their trained category set without extra fine-tuning. Open-vocabulary detection addresses this limitation. 

Instead of restricting \(\ell_i\) to a fixed training label set, open-vocabulary detectors compare image regions with text queries. This usually relies on a shared vision-language embedding space, where visual regions and textual labels can be scored by similarity. 

CLIP-style contrastive pretraining is a central foundation for this idea~\cite{pmlr-v139-radford21a}. Given a batch of matching image-text pairs, an image encoder maps each image to an embedding \(v_i\), and a text encoder maps each caption to an embedding \(t_i\). The training objective increases the similarity of matching pairs \((v_i,t_i)\) and decreases the similarity of non-matching pairs \((v_i,t_j)\), \(i\neq j\). 

A simplified image-to-text contrastive loss is 
\[ \mathcal{L}_{I\rightarrow T} = -\frac{1}{N}\sum_{i=1}^{N} \log \frac{\exp(\mathrm{sim}(v_i,t_i)/\tau)} {\sum_{j=1}^{N}\exp(\mathrm{sim}(v_i,t_j)/\tau)}, \] 
where \(\tau\) is a temperature parameter. A symmetric text-to-image term is usually added as well. 

After this training, textual labels can be compared directly with visual representations. OWL-ViT applies this principle to object detection by adapting a vision-language model so that image regions can be matched against text embeddings. The detector can therefore receive category names or ontology as text at inference time, rather than only predicting labels from a fixed classifier head~\cite{minderer_simple_2022}. 

OWLv2 extends this direction by scaling open-vocabulary detection with self-training: an existing detector generates pseudo-box annotations on large image-text data, and these pseudo-labels are then used to train a stronger detector~\cite{minderer_scaling_2024}. Open-vocabulary detectors are more flexible than closed-vocabulary detectors and they generally seem like a better choice, but one should notice their outputs depend on prompt wording, label granularity, and calibration between textual categories.
%Older pipelines exposed these steps as sliding windows, region proposals, classifiers, and post-processing. Modern detectors internalize much of the pipeline, but the functional needs remain: find objects, label them, and avoid returning many boxes for the same instance.
% TRACKING
Object detection is frame-local: it produces detections for one image, but it does not by itself say whether an object in frame \(t\) is the same physical instance as an object in frame \(t-1\). Multi-
  object tracking (MOT) addresses this problem by associating detections across time. In the common \emph{tracking-by-detection} formulation, a detector first produces a set of boxes for each frame,
  \[
  \mathcal{D}_t=\{(b_i^t,\ell_i^t,s_i^t)\}_{i=1}^{N_t},
  \]
  and the tracker maintains a set of tracks
  \[
  \mathcal{T}_t=\{\tau_j^t\}_{j=1}^{M_t}.
  \]
  Each track stores an estimated object state and history. At every time step, the tracker predicts where existing tracks should appear, compares these predictions with the new detections, and solves an
  assignment problem. A generic association step can be written as
  \[
  \hat{\pi}
  =
  \arg\min_{\pi}
  \sum_{(j,i)\in\pi}
  d(\tau_j^{t-1},d_i^t),
  \]
  where \(d(\tau_j^{t-1},d_i^t)\) is a matching cost between an existing track and a new detection. The Hungarian algorithm is commonly used to solve this assignment problem efficiently.

  SORT, short for \emph{Simple Online and Realtime Tracking}, is a representative example of this approach~\cite{bewley_simple_2016}. It combines a Kalman filter with Hungarian assignment. The Kalman
  filter predicts the next state of each tracked object, usually represented by bounding-box position, scale, aspect ratio, and their velocities. New detections are then matched to predicted tracks using
  bounding-box overlap, most commonly IoU. If a detection matches an existing track, the Kalman filter updates the track state. If a detection cannot be matched, a new track is created. If a track is not
  matched for several frames, it is removed.

  The strength of SORT is its simplicity and speed. It is detector-agnostic and can run very fast because it uses only motion prediction and box overlap. Its weakness follows directly from this design.
  SORT assumes that consecutive frames are close enough for a motion model and IoU overlap to remain meaningful. It has no appearance representation, so when two similar objects cross, overlap, or
  disappear behind occlusion, the tracker can switch their identities. It also depends heavily on detector quality: missed detections break tracks, and false detections may create spurious tracks.

  Deep SORT extends SORT by adding an appearance model~\cite{wojke_simple_2017}. For each detected crop, it extracts an appearance descriptor and combines this visual similarity with the motion-based
  Kalman prediction during assignment. This lets the tracker ask not only whether a detection is near the predicted position, but also whether it looks like the same object. The advantage is fewer
  identity switches when objects cross, overlap, or briefly disappear. The limitation is that the appearance model must be appropriate for the tracked domain, and the method still assumes a reasonably
  continuous video stream where motion prediction and recent appearance history remain meaningful.

  ByteTrack addresses a different weakness of tracking-by-detection pipelines~\cite{zhang_bytetrack_2022}. Many trackers discard low-confidence detections before association, which can break tracks
  during occlusion, blur, or poor lighting. ByteTrack first matches tracks to high-confidence detections and then tries to match still-unmatched tracks to low-confidence detections. The key idea is that
  a low detector score may still correspond to a real object. This improves robustness to temporary detector uncertainty, but it still relies on short-term spatial continuity and can be affected by false
  positives among low-confidence boxes.


  The common presupposition behind SORT, Deep SORT, and ByteTrack is therefore temporal and spatial continuity. Consecutive frames should be close in time, the camera viewpoint should not change too
  abruptly, and the same physical object should appear in nearby image locations across frames. Under these conditions, motion prediction, IoU overlap, and short-term appearance continuity are
  informative. When these conditions fail, tracking becomes much harder. The same object may reappear with no overlap with its previous box, a Kalman prediction may point to the wrong image region, and a
  short-term track may be deleted before the object is seen again.

  This motivates re-identification as a related but different problem. Re-identification does not try to estimate a continuous trajectory frame by frame. Instead, it asks whether two observations
  correspond to the same physical instance across different times, viewpoints, locations, or cameras. Zheng and Zheng define object re-identification as recognising object instances, not merely object
  categories, in different contexts of time, location, or viewpoint~\cite{zheng_object_2024}. This distinction is important: matching \enquote{the same chair} is different from detecting that both
  observations contain \enquote{a chair}.

  The standard Re-ID formulation is a query-gallery retrieval problem. Given a query crop \(x_q\), the system compares it against a gallery of previously observed crops
  \[
  \mathcal{G}=\{x_1,\ldots,x_n\}.
  \]
  A feature extractor maps each crop into an embedding space,
  \[
  f:x\mapsto f(x)\in\mathbb{R}^d.
  \]
  Similarity can then be computed with cosine similarity,
  \[
  \mathrm{sim}(x_q,x_g)=
  \frac{f(x_q)^\top f(x_g)}
  {\|f(x_q)\|\|f(x_g)\|},
  \]
  or with another learned metric. Ideally, gallery samples of the same instance should have higher similarity to the query than samples of different instances. In the person Re-ID literature, this is
  often formalized by splitting identities so that training and testing identities are disjoint, forcing the model to learn transferable identity-discriminative features rather than simply memorize
  training classes~\cite{gu_automated_2026}.

  Re-ID systems therefore focus on representation learning. The embedding should preserve instance-specific visual cues while being robust to nuisance variation such as viewpoint, illumination, scale,
  resolution, partial occlusion, and background change. Classical methods used hand-crafted descriptors and metric learning; modern methods usually use deep networks trained with identification losses,
  contrastive losses, triplet losses, or combinations of these objectives~\cite{zheng_object_2024}. The goal is not to classify an object into a category, but to place observations of the same instance
  close together in feature space.

  Re-identification is not a complete replacement for tracking. It does not provide a smooth trajectory, velocity estimate, or frame-by-frame state. However, it is a suitable substitute for the identity-
  association part when the assumptions of continuous tracking are violated. In a continuous video, SORT, Deep SORT, or ByteTrack can maintain identities using motion and short-term association. In
  sparse or discontinuous observations, appearance-based Re-ID is more appropriate because it can match an object after a long gap, a viewpoint change, or a large displacement in image coordinates. In
  practice, a robust system can combine both ideas: tracking handles short-term continuity, while Re-ID links observations across longer temporal and spatial gaps.


%TL;DR: Face recognition is useful for people but not a general persistence solution.
Pepper also provides robot-native people-perception modules through NAOqi, but these solve a narrower problem. ALPeoplePerception keeps track of people around the robot and exposes events and memory keys for detected, updated, or lost people. ALFaceDetection detects and can optionally recognize faces after a learning stage.\footnote{See the Aldebaran documentation for ALPeoplePerception and ALFaceDetection: \url{http://doc.aldebaran.com/2-4/naoqi/peopleperception/alpeopleperception.html} and \url{http://doc.aldebaran.com/2-4/naoqi/peopleperception/alfacedetection.html}.} These APIs are useful for human-centered interaction, but they are not general object-tracking or object-Re-ID systems. They primarily apply to people, depend on visibility and sensing conditions, and face recognition additionally depends on face pose, distance, and prior enrollment.

% %TL;DR: The section distinguishes the roles of detection, tracking, and Re-ID.
% The distinction among these methods is therefore functional. Detection answers what is visible and where it is. Tracking answers how identities persist in temporally continuous streams. Re-ID answers whether an instance seen now corresponds to one seen earlier under larger temporal or viewpoint changes. This distinction explains why scene-aware dialogue needs more than single-frame detection.

% Offloaded to Chapter~\ref{chap:server_side_rewrite}: the exact detector backends, open-vocabulary configuration, crop embedding model, thresholding, and robot-geometry consistency rules used by the server.

\section{Grounding and Scene Graph Generation}\label{sec:grounding-sgg}

%TL;DR: Grounding connects language to perceived entities and relations.
Chapter~\ref{chap:intro} introduced conversational grounding. In a robotic sense, grounding means connecting linguistic expressions to entities, attributes, and relations in the robot's perceived world. Lemaignan et al. describe situated grounding in HRI as a process that connects perception, symbolic knowledge representation, and dialogue processing~\cite[181--182]{lemaignan_grounding_2012}. The important point is that grounding is not merely adding an image to a prompt.

%TL;DR: Detections alone do not encode relations.
A list of detections is insufficient for many grounded questions. It may contain \((\text{cup}, b_1)\), \((\text{laptop}, b_2)\), and \((\text{table}, b_3)\), but it does not say whether the cup is on the table, left of the laptop, or currently being held by a person. Spatial and semantic dialogue often depends on exactly such relations.

%TL;DR: Scene graphs represent objects, attributes, and relations.
A scene graph represents a scene as a graph of objects and predicates. Following Johnson et al., scene graphs contain object nodes, object attributes, and relationships between objects~\cite[3668--3669]{johnson_image_2015}. A graph can be written as
\[
\mathcal{G}=(\mathcal{V},\mathcal{E}),
\]
where \(\mathcal{V}\) is a set of object instances and \(\mathcal{E}\) is a set of typed edges.

%TL;DR: Relations are usually represented as subject-predicate-object triplets.
A binary relation is often represented as a triplet
\[
(s,p,o),\qquad s,o\in\mathcal{V},\; p\in\mathcal{P},
\]
where \(s\) is the subject, \(p\) is the predicate, and \(o\) is the object. Attributes can be stored as node properties or as unary predicates, depending on the downstream representation. The general aim is the same: make visual facts queryable and compositional.

% Offloaded to Chapter~\ref{chap:server_side_rewrite}: the concrete attribute representation used by the server, including any self-relation encoding or schema normalization.

%TL;DR: SGG maps visual input to object and relation structure.
Scene Graph Generation (SGG) is the task of mapping visual input to object labels and relation triplets. If boxes and labels are given, the task focuses mainly on predicate prediction. If they are not given, SGG also includes object detection and classification. Zhu et al. survey SGG as a central task for high-level scene understanding and review methods based on feature extraction, relation fusion, and graph reasoning~\cite{zhu_scene_2022}.

%TL;DR: SGG is difficult because relations are numerous, biased, and ambiguous.
SGG is difficult for several reasons. The number of possible directed object pairs grows quadratically with the number of detected objects. Predicate distributions are long-tailed, so frequent relations such as \enquote{on} or \enquote{near} dominate rarer ones. Relation labels can also be ambiguous, because different annotators may choose different predicates for the same visual configuration.

%TL;DR: Detector errors propagate into relation errors.
SGG also depends strongly on object detection. If an object is missed, badly localized, or mislabeled, the correct relation may be impossible to recover downstream. This matters for dialogue because a wrong relation can become a wrong spoken claim about the user's visible environment.

%TL;DR: Rule-based relations are interpretable but semantically limited.
The simplest relation strategy is rule-based. Spatial predicates can be derived from geometry: left/right from projected centroids, above/below from vertical ordering, overlap from box intersections, and proximity from normalized distance. The advantage is interpretability and determinism. The disadvantage is semantic poverty: box geometry alone cannot reliably infer relations such as \enquote{holding}, \enquote{looking at}, or \enquote{using}.

%TL;DR: Learned SGG methods add semantic expressivity but inherit dataset limits.
Learned SGG models use visual features to predict relation triplets. They can model relations that geometry alone cannot capture, but they inherit dataset bias, label imbalance, and ontology limitations~\cite{zhu_scene_2022}. RelTR is one representative transformer-based method. It formulates relation prediction as set prediction with coupled subject-object queries, producing a sparse set of relation triplets directly~\cite{cong_reltr_2023}.

%TL;DR: VLM-based SGG improves openness but makes output control harder.
A newer alternative is VLM-based SGG. Instead of training a relation classifier for a fixed predicate set, a VLM can be prompted to generate relations from an image and a requested ontology. This is attractive for open or domain-specific environments. Work such as IndVisSGG shows how VLMs can support scene graph generation in industrial spatial intelligence~\cite{wang_indvissgg_2025}. The limitation is output control: a generative model may produce invalid structure, unsupported relations, or inconsistent wording.

%TL;DR: Strict output formats are useful but can constrain generative models.
Structured memory requires machine-readable output, but strict formats can also constrain generative reasoning. Tam et al. show that format restrictions such as JSON or XML can reduce LLM performance on reasoning and domain understanding tasks~\cite{tam_let_2024}. This does not mean that structure should be avoided. It means that architectures should separate what can be computed deterministically from what needs semantic interpretation.

%TL;DR: Set-of-Mark prompting makes visual references explicit.
Set-of-Mark prompting reduces referential ambiguity by marking visible objects with explicit identifiers before the image is passed to a VLM. The model can then refer to \enquote{object 3} rather than implicitly selecting among many similar regions. Yang et al. show that such marking can improve visual grounding behavior in GPT-4V-style systems~\cite{yang_set-of-mark_2023}. For scene-aware dialogue, this is relevant because object identifiers can connect visual evidence, graph nodes, and later language references.

%TL;DR: SGG metrics help, but dialogue needs more than single-image triplet recall.
SGG evaluation commonly distinguishes Predicate Classification (PredCls), Scene Graph Classification (SGCls), and Scene Graph Detection (SGDet)~\cite{zhu_scene_2022}. Metrics such as Recall@K and mean Recall@K measure whether predicted triplets match ground truth. These metrics are useful, but not sufficient for dialogue: a graph can have acceptable single-image triplet recall and still be unstable across turns or unhelpful for answering a user's question.

%TL;DR: Scene graphs are justified as structured context rather than as captions.
Captions and dense captions can describe an image in natural language, and systems such as the one proposed by Grassi et al. use dense captioning with LLMs to ground conversational robots on vision~\cite{grassi_grounding_2024}. Graphs have a different advantage. They are easier to update, query, compare, and constrain than free text. This makes them natural for targeted relational questions and persistent scene memory.

% Offloaded to Chapter~\ref{chap:server_side_rewrite}: the hybrid scene-graph construction strategy, deterministic spatial predicates, RelTR backend use, VLM prompting, Set-of-Mark rendering, schema validation, and graph merging.



\section{Grounding and Scene Graph Generation}\label{sec:grounding-sgg}

  %TL;DR: Grounding connects language to perceived entities and relations.
  Chapter~\ref{chap:intro} introduced conversational grounding. In a robotic sense, grounding means connecting linguistic expressions to entities, attributes, and relations in the robot's perceived
  world. Lemaignan et al. describe situated grounding in HRI as a process that connects perception, symbolic knowledge representation, and dialogue processing~\cite[181--182]{lemaignan_grounding_2012}.
  The important point is that grounding is not merely adding an image to a prompt.

  %TL;DR: Detections alone do not encode relations.
  A list of detections is insufficient for many grounded questions. It may contain \((\text{cup}, b_1)\), \((\text{laptop}, b_2)\), and \((\text{table}, b_3)\), but it does not say whether the cup is on
  the table, left of the laptop, or currently being held by a person. Spatial and semantic dialogue often depends on exactly such relations.

  %TL;DR: Scene graphs represent objects, attributes, and relations.
  Using standard graph-theoretic notation, a scene graph can be written as a tuple
  \[
  \mathcal{G}=(\mathcal{V},\mathcal{E},\lambda_{\mathcal{V}},\lambda_{\mathcal{E}}),
  \]
  where \(\mathcal{V}\) is a set of nodes, \(\mathcal{E}\subseteq \mathcal{V}\times\mathcal{V}\) is a set of directed edges, \(\lambda_{\mathcal{V}}\) assigns object labels or attributes to nodes, and
  \(\lambda_{\mathcal{E}}\) assigns relation labels to edges. In image-based scene graphs, nodes usually correspond to localized object instances and edges correspond to visual relations between them.
  Johnson et al.\ describe scene graphs as containing object instances, object attributes, and relationships between objects~\cite[3668--3669]{johnson_image_2015}. Zhu et al.\ similarly describe scene
  graphs as structured representations containing objects, attributes, and relationships, usually expressed through visual triplets~\cite{zhu_scene_2022}.

  %TL;DR: Relations are usually represented as subject-predicate-object triplets.
  A binary relation is usually represented as a triplet
  \[
  (s,p,o),\qquad s,o\in\mathcal{V},\; p\in\mathcal{P},
  \]
  where \(s\) is the subject node, \(p\) is the predicate, and \(o\) is the object node. For example, \((\text{cup},\text{on},\text{table})\) states a directed relation from the cup node to the table
  node. Attributes can be stored as node labels, node properties, or unary predicates such as \((\text{cup},\text{is},\text{red})\). The purpose is the same in each case: to convert visual information
  into a structured representation that can be queried and composed.

  % Offloaded to Chapter~\ref{chap:server_side_rewrite}: the concrete attribute representation used by the server, including any self-relation encoding or schema normalization.

  %TL;DR: SGG maps visual input to object and relation structure.
  Scene Graph Generation (SGG) is the task of automatically mapping an image or video into such a structured graph. It requires correct object localization, object classification, and relationship
  prediction~\cite{zhu_scene_2022}. If object boxes and labels are already given, the task reduces mainly to predicate prediction. If they are not given, SGG also includes detection and classification.

  Most SGG methods can be understood as either two-stage or one-stage methods. In a two-stage approach, objects are detected first, then pairs of detected objects are considered as candidate relations.
  For \(n\) detected objects, there are \(n(n-1)\) possible directed object pairs, so the number of candidate relations grows quadratically. A predicate classifier then predicts the relation label for
  each relevant pair. In a one-stage approach, the model attempts to predict objects and relations jointly, avoiding an explicit classification over every possible object pair~\cite{zhu_scene_2022}.

  The main difficulty of SGG lies in relationship prediction. Predicate distributions are strongly long-tailed: common relations such as \enquote{on}, \enquote{near}, or \enquote{has} occur far more
  often than rare predicates. As a result, models can become biased toward frequent generic relations. Relationship labels are also visually ambiguous. The same configuration may reasonably be described
  as \enquote{on}, \enquote{above}, \enquote{attached to}, or \enquote{part of}, depending on the annotation convention and task. Zhu et al.\ identify long-tailed distributions, intra-class diversity,
  contextual reasoning, and dataset bias as central problems in SGG~\cite{zhu_scene_2022}.

  %TL;DR: Detector errors propagate into relation errors.
  In two-stage SGG, object detection errors directly propagate into relation prediction. If an object is missed, no relation involving that object can be predicted. If a box is poorly localized, spatial
  predicates such as \enquote{on}, \enquote{inside}, or \enquote{next to} may become unreliable. If an object is mislabeled, the predicate model may predict a plausible relation for the wrong semantic
  pair. This dependency is one reason why SGG is harder than object detection alone.

  %TL;DR: Rule-based relations are interpretable but semantically limited.
  The simplest relation strategy is rule-based. Spatial predicates can be derived from bounding-box geometry: left/right from centroid order, above/below from vertical position, overlap from box
  intersection, and near/far from normalized distance. The advantage is interpretability and deterministic behavior. The limitation is semantic expressivity. Geometry can support simple spatial
  relations, but it cannot reliably infer relations such as \enquote{holding}, \enquote{wearing}, \enquote{using}, or \enquote{looking at}, which depend on visual context and object semantics.

  %TL;DR: Learned SGG methods add semantic expressivity but inherit dataset limits.
  Learned SGG models use visual features, object labels, object-pair features, context, and sometimes commonsense priors to predict relation triplets~\cite{zhu_scene_2022}. Two-stage methods commonly
  detect objects first, construct candidate subject-object pairs, extract features for each pair or their union region, and then classify the predicate. These methods are expressive, but they inherit the
  biases of their datasets, especially the long-tailed predicate distribution.

  RelTR is a representative one-stage transformer-based method. Inspired by DETR, it formulates SGG as set prediction over relation triplets rather than exhaustive classification of all object
  pairs~\cite{cong_reltr_2023}. The model uses coupled subject and object queries. A transformer encoder reasons over visual features, and a decoder predicts a fixed-size set of subject-predicate-object
  triplets. Training uses a matching loss between predicted and ground-truth triplets, encouraging sparse direct prediction of relations. This avoids explicitly enumerating every possible object pair,
  but the model still depends on the relation ontology and training distribution.

  %TL;DR: VLM-based SGG improves openness but makes output control harder.
  A newer alternative is VLM-based SGG. Instead of training a fixed predicate classifier, a vision-language model can be prompted to describe objects and relations in a requested schema or ontology. This
  is attractive when relation vocabularies are open, domain-specific, or not well covered by standard SGG datasets. Work such as IndVisSGG shows how VLMs can be used for scene graph generation in
  industrial spatial-intelligence settings~\cite{wang_indvissgg_2025}. The limitation is output control: a generative model may produce unsupported relations, inconsistent predicate names, invalid
  identifiers, or malformed structure. For that reason, VLM-based SGG often needs validation, normalization, and filtering after generation.

  % MOVE TO SECTION~\ref{sec:llms-vlms}: strict structured-output formats and Set-of-Mark prompting are primarily VLM prompting issues rather than core SGG definitions.
  %
  % Structured memory requires machine-readable output, but strict formats can also constrain generative reasoning. Tam et al. show that format restrictions such as JSON or XML can reduce LLM performance
  on reasoning and domain understanding tasks~\cite{tam_let_2024}. This does not mean that structure should be avoided. It means that architectures should separate what can be computed deterministically
  from what needs semantic interpretation.
  %
  % Set-of-Mark prompting reduces referential ambiguity by marking visible objects with explicit identifiers before the image is passed to a VLM. The model can then refer to \enquote{object 3} rather
  than implicitly selecting among many similar regions. Yang et al. show that such marking can improve visual grounding behavior in GPT-4V-style systems~\cite{yang_set-of-mark_2023}. For scene-aware
  dialogue, this is relevant because object identifiers can connect visual evidence, graph nodes, and later language references.

  %TL;DR: SGG metrics help, but they must be interpreted carefully.
  SGG evaluation commonly distinguishes three settings~\cite{zhu_scene_2022}. In Predicate Classification (PredCls), ground-truth boxes and object labels are given, and the model predicts only
  predicates. In Scene Graph Classification (SGCls), ground-truth boxes are given, but object labels and predicates must be predicted. In Scene Graph Detection (SGDet), boxes, object labels, and
  predicates must all be predicted. These settings progressively increase difficulty.

  A common metric is Recall@K. For each image, the model ranks predicted triplets by confidence. Recall@K measures the fraction of ground-truth triplets that appear among the top \(K\) predictions:
  \[
  \mathrm{Recall@K}
  =
  \frac{|\mathcal{T}_{gt}\cap \mathcal{T}_{pred}^{K}|}
  {|\mathcal{T}_{gt}|},
  \]
  where \(\mathcal{T}_{gt}\) is the set of ground-truth triplets and \(\mathcal{T}_{pred}^{K}\) is the set of top-\(K\) predicted triplets. Because ordinary Recall@K is dominated by frequent predicates,
  mean Recall@K computes recall separately for each predicate class and then averages across classes. This gives rare predicates more influence and is often more informative under long-tailed predicate
  distributions~\cite{zhu_scene_2022}.

  %TL;DR: Scene graphs are justified as structured context rather than as captions.
  Captions and dense captions can describe an image in natural language, and systems such as the one proposed by Grassi et al.\ use dense captioning with LLMs to ground conversational robots on
  vision~\cite{grassi_grounding_2024}. Scene graphs serve a different role. They are structured, typed, and easier to update, query, compare, and constrain than free text. This makes them useful when
  downstream components need explicit objects, attributes, and relations rather than a single natural-language description.

  % Offloaded to Chapter~\ref{chap:server_side_rewrite}: the hybrid scene-graph construction strategy, deterministic spatial predicates, RelTR backend use, VLM prompting, Set-of-Mark rendering, schema
  validation, and graph merging.


\section{Language Models and Vision-Language Models}\label{sec:llms-vlms}

%TL;DR: Language modeling is probabilistic sequence modeling.
Natural language processing (NLP) is the computational study and processing of human language. Language modeling is a narrower task inside NLP: estimating probabilities over token sequences. For a token sequence \(x_{1:T}\), an autoregressive language model factorizes
\[
p(x_{1:T})=\prod_{t=1}^{T}p(x_t\mid x_{<t}).
\]
Generation is repeated next-token prediction. This definition is important because a language model is optimized to produce text under context, not to be truthful by construction.

%TL;DR: Earlier NLP methods were explicit but limited in context.
Before the current Transformer-based paradigm, language modeling included n-gram models, naive Bayes classifiers for many classification tasks, and later recurrent neural networks. N-gram models are interpretable and efficient, but they suffer from sparsity and limited context. Recurrent neural networks improve context modeling through hidden states, and gated variants such as LSTMs and GRUs handle longer dependencies better than plain recurrence. However, recurrent computation remains sequential, which limits large-scale parallel training.

%TL;DR: Transformers made large-scale pretraining practical.
The Transformer architecture replaces recurrence with self-attention~\cite{vaswani2023attentionneed}. Self-attention lets tokens attend to other tokens in the context and allows much more parallelism during training than recurrent models. Large language models (LLMs) are usually large Transformer-based models trained on broad corpora and then adapted for instruction following or dialogue.

%TL;DR: Instruction tuning improves usability but does not provide grounding by itself.
A base pretrained model is not automatically a dialogue agent. It may complete text without following the user's intent. Instruction tuning and alignment methods modify this behavior using supervised instruction data and preference-based training. InstructGPT-style training combines supervised fine-tuning with reinforcement learning from human feedback~\cite{ouyang2022traininglanguagemodelsfollow}. This improves interaction, but it does not remove the need for external grounding.

%TL;DR: Decoding policy affects whether answers are stable or variable.
The generated answer also depends on decoding. Greedy decoding chooses the most likely next token, beam search keeps several high-probability continuations, and stochastic decoding samples from a probability distribution. Temperature controls how deterministic or varied the output is. For robotics, this is not only a stylistic parameter: grounded answers often need stability more than creative variation.

%TL;DR: VLMs condition language on visual input.
Vision-language models (VLMs) extend language models with visual input. Architectures differ in how they compute and fuse visual representations. Some use contrastive dual encoders for retrieval or classification. Others project visual tokens into the LLM input space and train the model to follow multimodal instructions. Examples include LLaVA, InstructBLIP, and PaLM-E~\cite{liu_visual_2023,dai_instructblip_2023,driess_palm-e_2023}.

%TL;DR: Contrastive image-text learning supports open-vocabulary behavior.
Contrastive vision-language pretraining is central to open-vocabulary perception. CLIP-style training aligns images and texts in a shared embedding space by bringing matching image-text pairs closer and pushing mismatched pairs apart~\cite{pmlr-v139-radford21a}. This shared geometry supports zero-shot classification, open-vocabulary detection, retrieval, and prompt-conditioned visual reasoning.

%TL;DR: VLMs are semantically powerful but can hallucinate.
VLMs are useful because they connect visual evidence with linguistic categories and relations. Their failure modes are also important. They can hallucinate visually unsupported objects or relations, describe similar objects inconsistently across turns, or fail to produce strict schemas. In embodied dialogue, such errors are visible to the user and therefore become grounding errors, not merely language errors.

%TL;DR: Robotics LLM/VLM work often focuses on planning rather than persistent dialogue memory.
Several robotics systems use LLMs or VLMs to connect language with action. SayCan grounds high-level language model suggestions through affordance scores from robot skills~\cite{ahn_as_2022}. Code as Policies uses language models to generate robot policy code from natural language instructions~\cite{liang_code_2023}. PaLM-E conditions multimodal language modeling on embodied inputs~\cite{driess_palm-e_2023}. These systems are important relatives, but their central question is often action selection or planning rather than persistent scene dialogue.

%TL;DR: For scene-aware dialogue, language models are conditional generators over external context.
For the present problem, the most useful view is that LLMs and VLMs are conditional generators. The question is not whether the model independently \enquote{knows} what is in the room. The question is whether the system supplies reliable visual and structured context. This leads directly to retrieval and caching, because scene memory must be converted into usable context for generation.

% Offloaded to Chapter~\ref{chap:server_side_rewrite}: exact LLM/VLM provider choices, prompts, decoding parameters, structured-output repair, and dialogue-service behavior.

\section{RAG, CAG, and GraphRAG}\label{sec:rag-cag-graphrag}

%TL;DR: RAG separates language generation from external evidence access.
Retrieval-Augmented Generation (RAG) combines a generator with external evidence retrieval. Lewis et al. describe RAG as a combination of parametric memory in a pretrained model and non-parametric memory in a retrievable corpus~\cite{lewis_retrieval-augmented_2020}. Given a query, a retriever selects relevant context, and the generator produces an answer conditioned on both the query and that context.

%TL;DR: Retrieval quality introduces precision-latency tradeoffs.
In many RAG systems, retrieval is embedding-based. A query and document chunks are mapped into a vector space, and relevance is approximated by cosine similarity or inner product. Approximate nearest-neighbor indexing makes retrieval efficient for large corpora, but it introduces tradeoffs between speed, recall, embedding dimensionality, chunking, and reranking. In interaction, these tradeoffs affect both answer quality and response latency.

%TL;DR: RAG is useful for changing knowledge, but naive chunking can destroy structure.
RAG is useful when the model needs external, changing, or too-large knowledge. In scene-aware dialogue, the external knowledge is not a normal document collection but the robot's current and past scene state. A naive textual RAG approach could serialize the scene graph into chunks, but chunking can destroy relations. If \enquote{cup}, \enquote{table}, and \enquote{on} are separated, the generator receives fragments rather than a scene fact.

%TL;DR: Multimodal RAG generalizes retrieval beyond text.
Recent retrieval-augmented work also extends beyond text into images and multimodal data~\cite{zheng_retrieval_2025}. Older multimodal RAG pipelines often converted images to captions and retrieved text. Newer systems can retrieve across text, image, table, or other representations using multimodal embeddings. This is relevant background, although scene-aware dialogue often needs structured retrieval over entities and relations rather than only semantic similarity.

%TL;DR: CAG preloads bounded knowledge into context.
Cache-Augmented Generation (CAG) uses a different assumption. Instead of retrieving external chunks at every turn, it preloads relevant knowledge into the model context or cache and reuses it during generation. Chan et al. argue that, for some bounded knowledge tasks, CAG can be sufficient and can avoid RAG's online retrieval overhead~\cite{Chan_2025}. The prerequisite is that the relevant knowledge fits into the usable context and does not change too rapidly.

%TL;DR: CAG moves selection from query time to cache construction time.
CAG does not remove selection; it moves selection earlier. In RAG, the system decides what to retrieve after the user asks a question. In CAG, the system decides what to include when constructing the cache. This is useful for broad scene awareness, but less useful when a question requires precise retrieval around one object or relation.

%TL;DR: GraphRAG retrieves through relational structure.
GraphRAG introduces explicit relational structure into retrieval. Instead of treating memory as independent chunks, it represents memory as a graph and retrieves evidence through traversal, neighborhood expansion, or graph summaries. This is natural when the source representation is already relational. A question about an object can retrieve its local neighborhood rather than an unrelated textual chunk.

%TL;DR: Scene graphs connect naturally to graph-based retrieval.
Open-world 3D scene graph work has begun to connect scene graphs with retrieval-augmented reasoning~\cite{yu_open-world_2025}. VLM-based graph representations for robotics and industry similarly show that graph structure can support downstream reasoning and task execution~\cite{li_vlm-msgraph_2025,wang_indvissgg_2025}. The important point for this thesis is conceptual: the evidence-access mechanism should match the structure of the knowledge state.

%TL;DR: Retrieval mechanisms must be evaluated as interaction mechanisms.
For social robots, retrieval quality is not only a question-answering metric. A slow but accurate answer may break conversational rhythm. A fast answer that contradicts scene memory may damage trust. A broad cache may support scene summaries but fail on specific relations. A graph query may answer precise questions but require reliable entity matching first.

% Offloaded to Chapters~\ref{chap:system_architecture_rewrite} and~\ref{chap:server_side_rewrite}: the concrete split between cached scene context, object-focused graph retrieval, dialogue context assembly, and any query-routing heuristics.

\section{Pepper Robot as an Embodied Platform}\label{sec:pepper-platform}

%TL;DR: Pepper is introduced as the concrete platform that imposes constraints.
Pepper is a humanoid social robot platform designed for proximate human interaction. It has speech, speakers, microphones, cameras, a tablet, expressive motion, and a mobile base. For this thesis, Pepper is best understood as an embodied communication platform rather than as a manipulation platform. Its main output is not physical manipulation, but situated interaction about the visible environment.

%TL;DR: Pepper's hardware is limited relative to modern AI workloads.
Pepper's onboard hardware is limited relative to modern vision-language workloads. Public technical specifications list an Intel Atom E3845 CPU, 4 GB RAM, and 32 GB eMMC flash memory in the main CPU module, together with cameras, microphones, a tablet, sonars, lasers, and other sensors.\footnote{See the Pepper technical specifications: \url{https://aldebaran.com/support/kb/pepper/general/pepper-technical-specifications/}.} These resources are sufficient for native robot control and interaction modules, but not for running a full contemporary perception-to-language pipeline onboard.

%TL;DR: NAOqi exposes robot functions but does not provide open-ended grounded dialogue.
Software-wise, Pepper 2.5 uses NAOqi, a middleware that exposes robot functions through modules and APIs~\cite{softbank_python_sdk_2026}. NAOqi provides access to sensing, speech, tablet output, behavior control, and predefined speech-recognition functionality~\cite{softbank_alspeechrecognition_2026}. Such grammar- or phrase-based interaction is useful for bounded commands, but it does not by itself provide open-ended grounded dialogue.

%TL;DR: Pepper's tablet is part of the interaction channel.
Pepper's tablet is not merely an attached display. It is an interaction modality that can show options, present visual state, or support disambiguation~\cite{softbank_peppers_tablet_2026}. In grounded dialogue, this matters because visible feedback can reduce the mismatch between the robot's internal state and the user's interpretation.

%TL;DR: Existing Pepper systems commonly rely on external computation.
Several related Pepper systems already use external processing or modern AI components. Reyes et al. demonstrate near real-time object recognition for Pepper using deep neural networks running outside the robot~\cite{reyes_near_2018}. Trinquet et al. combine Pepper with object detection and OCR for library assistance~\cite{trinquet_future_2025}. Brad et al. propose a RAG architecture for Pepper in industrial assistance~\cite{brad_retrieval-augmented_2025}. Latif et al. build an LLM-powered interactive learning robot with YOLO-based perception~\cite{latif_physicsassistant_2024}.

%TL;DR: The platform motivates distributed scene-aware dialogue.
These works show that distributed architectures are a recurring response to Pepper's computational limits. They also show that Pepper is a useful platform for studying grounded interaction because it forces the full stack to be considered: sensing, dialogue, latency, embodiment, user expectations, and external computation. The specific implementation choices are therefore described later as architectural responses to these platform constraints.

% Offloaded to Chapters~\ref{chap:robot_side_rewrite} and~\ref{chap:server_side_rewrite}: exact robot-side services, image capture, metadata collection, tablet behavior, server communication, and runtime orchestration.

% \section{Summary}\label{sec:background-summary}

% %TL;DR: The chapter connected HRI, perception, scene representation, language, retrieval, and Pepper.
% This chapter introduced the background needed for the rest of the thesis. HRI and social robotics explain why Pepper's embodiment makes interaction quality part of the technical problem. The problem setting then framed scene-aware dialogue as dialogue over persistent scene memory rather than over a single image.

% %TL;DR: The perception sections explained evidence and persistence.
% Object detection provides localized visual evidence, but closed-vocabulary detectors are limited by their ontology and open-vocabulary detectors introduce prompt and calibration sensitivity. Tracking adds temporal identity, but many tracking methods assume continuity that discrete robot scans can violate. Re-identification is therefore the relevant concept for larger viewpoint and temporal gaps.

% %TL;DR: Scene graphs were introduced as the bridge between perception and dialogue.
% Grounding requires more than detections or captions. It requires a representation in which objects, attributes, and relations can be queried and updated. Scene graphs provide such a representation, while SGG methods differ in their tradeoffs between interpretability, semantic expressivity, openness, and output control.

% %TL;DR: Language and retrieval methods explain how memory becomes dialogue context.
% LLMs and VLMs provide language and multimodal interpretation, but they do not remove the need for external world state. RAG, CAG, and GraphRAG describe different ways of making external state available to generation. The correct choice depends on the structure, size, and update frequency of the knowledge state.

% %TL;DR: The method chapters now have a clean theoretical basis.
% The main result of this chapter is a set of assumptions and terms. The system observes scenes discontinuously, must preserve object identity, must represent relations, must answer in natural language, and must offload heavy computation away from Pepper. The following chapters describe how the implemented architecture satisfies these requirements and where the resulting tradeoffs appear in practice.
